\section{From complete binary/ternary C to actual coefficient unimodality}
This is a theorem about the ORIGINAL coefficients, with no amplification.
Let $\mathcal B_3$ be the normalized increasing arms
$1^\alpha2^u3^r$, $\alpha\ge1$, $u,r\ge0$, together with the empty
arm when a deletion reaches it. An initial part of size $a$ can be
replaced by $a$ ones, so this normalization preserves arm polynomials.
All coefficients are zero extended. Write $[t]=1+\cdots+z^{t-1}$,
$S_t=z[t]$, $q_X=z^{a}F_X(1/z)$, $a=\operatorname{wt}X$.

The proof gives a larger class $\mathcal A_4$: the union of
$\mathcal B_3$, all $1^\alpha2^uT$ ($T\ge2$), and all arms having
at most TWO arbitrary non-one parts. Include rectangles and the empty
arm. Every member has complete C: use the new binary/ternary theorem
or the existing at-most-two-nonbinary-capacities theorem in
\texttt{two\_nonbinary\_complete\_C\_proof.tex}.
Every nonrectangular member has $3\le\kappa\le4$.
For binary-prefix one-jump arms use $\kappa(1^\alpha2^u)=3$ when
$u>0$, and $\kappa(ZT)=\min(T+1,\kappa(Z)+1)$; for two-part
skeletons use $\kappa\le h+2$. Both facts include arbitrary padding.
This union is closed under final-part deletion, and its pair family
is closed under arm exchange.

\begin{theorem}[Full A2 and full source shape for all binary/ternary pairs]
For every $X,Y\in\mathcal A_4$, put $p=q_Xq_Y$ and
$d_n=(p_n-p_{n-1})^+$. Then $(2^nd_n)_{n\ge0}$ is unimodal.
If both arms are nonrectangular, the positive part of the ENTIRE
raw source $E=(2-z)p$ is unimodal, its positive support is an initial
interval, and all its other coefficients are nonpositive.
Both arm counts and initial runs of ones are unrestricted.
\end{theorem}

\subsection{An unconditional tail resource independent of A2}
The complete C theorems for the specified arms in $\mathcal A_4$
are the single-arm inputs; the only new one is the all-count
binary/ternary theorem. They are not finite diagnostics. The clipped-ratio argument in
\texttt{clipped\_forward\_ratio\_frontier.tex} gives coefficient tilt A1:
$2^np_n$ is unimodal for every nonrectangular pair in this class.
Consequently $E_i<0$ implies $E_j\le0$ at every $j>i$.
This consequence does NOT assume A2.

We next prove, independently of A2, that if $\nu$ is the LAST
ordinary mode of a nonrectangular pair, then
\[
 n>\nu,\quad E_n\ge0\quad\Longrightarrow\quad E_{n+1}\le E_n.
 \tag{Q}
\]
Put $D=a+b$ and $j=D-n$. At $j\ge2$ classify each ordinary forward
arm by whether $F(j+1)>F(j)$.
If both rise, complete C and the source-prefix theorem in
\texttt{C\_to\_coupled\_S\_rising\_window.tex} prove (Q), including
the case where the source at $n$ is zero. If both do not rise,
the first-half theorem in \texttt{agent\_bothdeclining\_first\_birth.tex}
applies when $2j\le D+2$. Otherwise $2n\le D$, and the actual
post-LAST-mode theorem in
\texttt{agent\_half\_lc\_actual\_weak\_curvature.tex} applies.
Neither of these two nonrising resources needs an A2 hypothesis.

Here is a direct proof of the mixed case that also includes arms
with only ONE non-one part. Orient $A=F_Z$ nonincreasing at $j+1$
and $B=F_Y$ strictly rising there. Its source
$E_Y=1-z+(2z-1)B$ has nonnegative increments through $j+1$.
Its first positive increment is at $3\le\kappa\le4$, of size at least
one, because $Y\in\mathcal A_4$ is nonrectangular.
Put $Q=E_YA$, $r=a-j$, and $s_r=\Delta q_Z(r)$. The exact identity is
\[
 E_{n+1}-E_n=(s_{r+1}-s_r)-(Q_{j+1}-Q_j),\qquad
 Q_{j+1}-Q_j=\sum_i\Delta E_Y(i) A_{j+1-i}.       \tag{1}
\]
The sum is nonnegative. For $r\le-2$ the first term vanishes.
For $r=-1$ it equals one, and its compulsory birth term has
index $a+2-\kappa\in[1,a-1]$ since $a\ge3$; hence that term pays it.

For $r\ge0$, a putative positive value in (1) forces
$d_{r+1}>d_r>0$ for the owner's actual reverse positive differences.
Let $H=\sum_{t\in Z,\ t>1}(t-1)$, $M=\alpha+H$, and
$C=d_{r+1}-d_r>0$. The established single-arm skeleton bound is
\[
 H\ge2r+2,\qquad (r+1)C\le(H-2r-1)d_r.           \tag{2}
\]
Every positive box of $q_Z$ has endpoint sum at least $M$.
Thus $q_Z(v)\ge q_Z(w)$ if $w\le v$ and $w+v\le M$.
If $2r+\kappa\le M$, take $w=r+1$, $v=r+\kappa-1$;
then the compulsory term $A_{j+1-\kappa}=q_Z(v)$ pays
$C\le q_Z(r+1)$.
Otherwise $\kappa\le4$, $\alpha\ge1$, and $H\ge2r+2$
force EXACTLY $\kappa=4$, $\alpha=1$, $H=2r+2$.
Now $r+v=M$, so $q_Z(v)\ge q_Z(r)\ge d_r\ge C$ by (2).
These reflection comparisons also prove that the birth is in the
actual support, so no absent source is charged. This proves the
mixed case, without a finite entrance or a minimum count of two.

At the remaining reverse terminal ranks, $j=1$ gives
$E_n=2(AB)_1-(AB)_2=4-5=-1$, so its nonnegative premise fails.
At $j=0$, $E_n=0$, $E_{n+1}=-1$; the next rank has negative source,
and beyond it the source is zero. Negative forward ranks are zero.
Thus (Q) holds at EVERY required rank.

\subsection{Full old A2 supplies a signed source suitable for every window}
Suppose A2 holds for a nonrectangular pair in $\mathcal A_4$.
Let $\mu,\nu$ be the first and last ordinary modes of $p$.
Through $\mu$, $d_n=\Delta p_n$ and the exact identity
\[
 2^{n+1}d_{n+1}-2^nd_n=2^n(E_{n+1}-E_n)          \tag{3}
\]
transfers A2 to unimodality of $E_0,\ldots,E_\mu$.
All these source values are positive. Over the maximum plateau,
$E_\mu\ge\max p$ and $E_{\mu+1}=\cdots=E_\nu=\max p$;
also $E_{\nu+1}\le\max p$. For later comparisons (Q) pays whenever
the current source is nonnegative; if it is negative, A1 precludes
any later positive source. It follows that $E^+$ is unimodal with
initial positive support. Zero values after its mode cannot be
followed by a positive value, by (Q); negative values cannot be
followed by a positive value, by A1. This checks zeros and the
ENTIRE plateau, rather than identifying first and last modes.

The signed-window lemma of \texttt{signed\_kernel\_extension.tex}
therefore makes $([t]E)^+$ unimodal with initial positive support
for EVERY integer $t\ge1$. The lemma applies to the whole raw
source, including its negative terminal coefficient.

\subsection{A full A2 induction, including its bases}
First suppose one arm is rectangular, $q_X=[\alpha+1]$, and put
$b=\deg q_Y$. If $\alpha\ge b$, the first ordinary product mode is
$b$. Its positive differences are $(q_Y)_0,\ldots,(q_Y)_b$, then
zeros. Their weighted values increase, since actual ideal growth
gives $2(q_Y)_n\ge(q_Y)_{n-1}$. If $\alpha<b$, the first product
mode $\sigma$ is below $b$: at rank $b$ the window difference is
$1-(q_Y)_{b-\alpha-1}\le0$, and ordinary window unimodality applies.
The single-arm corrected source
$e_Y=(2-z)q_Y-z^b+z^{b+1}$ is nonnegative unimodal.
The product source equals $[\alpha+1]e_Y$ at every rank below $b$.
Its prefix through $\sigma$ is unimodal, and (3), followed by the
zero positive-difference tail, proves A2. This includes the empty
arm and two rectangles. We do NOT assert full raw-source shape
for rectangular pairs.

Induct on the total number of non-one parts in a nonrectangular
$\mathcal A_4$ pair and orient a heavier arm as $X$.
If $X$ has only one non-one part, the full heavier-one-part A2
theorem of \texttt{deficit\_cumulative\_ratios.tex} applies.
Otherwise write $X=Zt$ with its actual legal final capacity $t$;
both $Z,Y$ are still
nonrectangular and belong to $\mathcal A_4$. Full OLD A2 holds by
the smaller-count hypothesis. The preceding source argument applies
to $p=q_Zq_Y$. The exact NEW product and source are
\[
 q_Xq_Y=[t]p+z^a q_Y,\qquad
 (2-z)q_Xq_Y=[t]E+z^a(2-z)q_Y.                 \tag{4}
\]
The NEW first ordinary mode $\sigma$ satisfies $\sigma<a$,
including equal weights, by the strengthened ordinary midpoint
theorem. Hence its source through $\sigma$ equals $[t]E$.
These source values are positive, so they equal $([t]E)^+$.
Signed-window unimodality and (3) prove NEW A2 through $\sigma$;
all later NEW positive differences vanish. The induction closes.
Applying the preceding source argument once more to the now-proved
full A2 gives the full-source assertion of the theorem.
Only FULL old A2 is used here; a truncated old hypothesis is not
substituted at any window read.

\begin{corollary}[An arbitrary heavier final jump]
Let $Z,Y\in\mathcal A_4$, append ANY legal final part $t$ to $Z$,
and assume $a=\operatorname{wt}(Zt)\ge\operatorname{wt}Y$.
Then A2 holds for $q_{Zt}q_Y$, with no upper bound on $t$.
\end{corollary}
\begin{proof}
For both old arms nonrectangular, (4), full old A2, A1, (Q), and
the signed-window lemma apply for this arbitrary $t$ as well.
If $Z$ is rectangular, use the heavier-one-part base; if $Y$ is
rectangular, use the rectangular base. Empty bases are included.
\end{proof}

\begin{theorem}[Original coefficient theorem with an arbitrary strict peak]
Let $X,Y\in\mathcal A_4$ and let $M>\max(X,Y)$ be any strict peak.
Then the ORIGINAL core and BOTH relative polynomials
\[
 F_X+F_Y+(S_M-1)F_XF_Y,\quad
 F_X(1+(S_M-1)F_Y),\quad F_Y(1+(S_M-1)F_X)
\]
are coefficient unimodal. The same holds if one arm is $Zt$,
$Z,Y\in\mathcal A_4$, for ONE arbitrary legal final jump $t$,
provided $\operatorname{wt}(Zt)\ge\operatorname{wt}Y$.
The strict peak $M$ is unrestricted above the arm capacities.
\end{theorem}
\begin{proof}
Use the least-peak, least-relative-counterexample route in
\texttt{inductive\_signed\_window.tex}, or its shorter-prefix
refinement in \texttt{short\_prefix\_reflection\_join.tex}.
The family of all $\mathcal A_4$ pairs, adjoined to the stated
heavier-final-jump pairs, is closed under arm exchange and deleting
a final part. Deleting the exceptional jump gives a $\mathcal A_4$
pair; deleting from the opposite arm preserves its weight inequality.
The full A2 statements just proved supply the required weighted
positive-difference input, while deletion of a proper relative
polynomial supplies the signed prefix by leastness.

For completeness the lower-peak base also stays in this family.
Let $t_0$ be the largest arm capacity. For an $\mathcal A_4$
pair, removing all occurrences of $t_0$ leaves another such pair.
For a qualifying exceptional jump, if $t_0=t$, remove this jump
first and obtain an $\mathcal A_4$ pair. If $t_0>t$, all occurrences
of $t_0$ lie in the opposite arm: removing them preserves the
qualifying heavier arm. Continue stripping all maxima to obtain a
strict peak $t_0<M$ in the same family. Thus least-peak induction
or the rectangular base applies.
Restore each removed maximum by $J=1+S_{t_0}J^{\rm smaller}$.
The existing base-peak join contradicts the balanced relative
counterexample. Both relative assertions follow, and heavier-first
core-relative equivalence supplies the core assertion everywhere.
\end{proof}

\paragraph{Scope and verification.}
The same proof applies to ANY final-deletion-closed class of normalized
arms with complete C and first missing weight at most four; the stated
$\mathcal A_4$ is an explicit class for which all these inputs are proved.
These are unbounded original COEFFICIENT theorems, not inferences
from cumulative log-concavity and not conclusions of finite tests.
They retain arbitrarily large peaks, arbitrarily many twos and
threes, and the stated arbitrary heavier jump. They do NOT solve
the original conjecture for arbitrary arm capacities or general C/A2.
The new Lean file checks the small-first-missing mixed payment,
whole-function positive-source splice and absence of holes, exact
source/weighted-difference identity, and the relevant mode/source
index arithmetic. The actual arm membership, prior shape resources,
signed-window and least-counterexample assemblies are written proofs;
no end-to-end Lean formalization of compositions is claimed.
